\notepage{32}{The Heisenberg Equation of Motion}
For a unitary transformation $U$, there are two equivalent ways to
calculate a transformed matrix element:
\begin{center}
\begin{tabular}{lll}
\toprule
 & Transform the states & Transform the operator\\
\midrule
Ket & $\ket\psi\mapsto U\ket\psi$ & $\ket\psi$ fixed\\
Bra & $\bra\phi\mapsto\bra\phi U^\dagger$ & $\bra\phi$ fixed\\
Operator & $\Omega$ fixed & $\Omega\mapsto U^\dagger\Omega U$\\
\bottomrule
\end{tabular}
\end{center}
Both prescriptions give
\[
 (\bra\phi U^\dagger)\Omega(U\ket\psi)
 =\mel\phi{U^\dagger\Omega U}\psi.
\]
Either can be used according to convenience. For example, parity gives
$(\bra{x'}\Pi^\dagger)X(\Pi\ket x)=\mel{x'}{\Pi^\dagger X\Pi}x$.
\smallheading{Two pictures of time evolution}
In the Schr\"odinger picture,
$\ket{\psi(t)}=U(t)\ket{\psi(0)}$, while an operator without explicit
time dependence remains fixed. In the Heisenberg picture the state is
fixed at $\ket{\psi(0)}$ and the operator evolves:
\[
 \Omega_H(t)=U^\dagger(t)\Omega U(t).
\]
The Schr\"odinger equation for arbitrary initial states implies
\[
 \ii\hbar\frac{\dd}{\dd t}U(t)\ket{\psi(0)}
 =HU(t)\ket{\psi(0)},\qquad
 \ii\hbar\frac{\dd U}{\dd t}=HU.
\]
Differentiating $U^\dagger\Omega U$ gives
\[
 \boxed{\ii\hbar\frac{\dd\Omega_H}{\dd t}
 =\comm{\Omega_H}{H_H}.}
\]
This is useful when only expectations, rather than the full wavefunction,
are needed. An explicitly time-dependent Schr\"odinger operator adds
$\ii\hbar\,U^\dagger(\partial\Omega/\partial t)U$ to the right side.
\smallheading{Comparison with Hamilton's equations}
For $H=P^2/(2m)+V(X)$ and $\comm XP=\ii\hbar I$,
\begin{align*}
 \dot x&=\frac{\partial\mathcal H}{\partial p}=\frac pm,
 &\dot X&=\frac1{\ii\hbar}\comm XH=\frac Pm,\\
 \dot p&=-\frac{\partial\mathcal H}{\partial x}=-V'(x),
 &\dot P&=\frac1{\ii\hbar}\comm PH=-V'(X).
\end{align*}
The commutator identities used here are
$\comm X{f(P)}=\ii\hbar f'(P)$ and $\comm P{g(X)}=-\ii\hbar g'(X)$,
for functions and domains where these formal derivatives are valid.
